\chapter*{ВСТУП}
\addcontentsline{toc}{chapter}{ВСТУП}

Сучасні комп'ютери, і особливо їх графічні процесори, можуть обробляти все більше даних за більш короткий, що відкриває нові можливості в галузі комп'ютерної графіки. Однак, більшість відеоігор та інших графічних застосунків досі використовує більш продуктивний, але менш фізично-коректний і реалістичний метод рендерингу --- растеризацію, замість більш ресурсно-затратної технології --- трасування променів.

Растеризація --- це процес трансформації векторних даних (зазвичай множини трикутників) в растрове зображення, яке графічний процесор по пікселю відмальовує на екрані. В той же час трасування променів \engl{ray tracing} --- це зовсім інший підхід, який полягає в отриманні інформації про об'єкти рендерингу через кидання променю з кожного пікселя на їх поверхню. Із самого означення зрозуміло, що наївна реалізація алгоритму трасування променів досить ресурсно-затратна, і тому вимагає оптимізацій для того, щоб можна було використати цей алгоритм на практиці.

Метою цієї роботи є розглянути основні поняття, що стосуються трасування променів, в тому числі воксельного, різні методи оптимізації рей трейсингу, їх поєднання з методом трасування у воксельному просторі, а також переваги та недоліки цього методу.

Дана робота ставить на меті виконання наступних завдань:

\begin{itemize}[itemsep=0pt, topsep=0pt]
    \item Аналіз та порівняння різних методів трасування у дискретному та неперервному просторі;
    \item Реалізація алгоритму трасування воксельних об'ємів і створення проміжних даних для подальшого рендерингу;
    \item Реалізація пост-процесингового алгоритму відкладеного \engl{defer\-red} трасування у воксельному просторі з використанням загального воксельного об'єму, створеного на основі інших об'єктів 3D-сцени;
    \item Дослідження додаткових методів оптимізації трасування воксельної сітки;
    \item Створення кінцевого застосунку, що демонструє всі наведені вище алгоритми рендерингу на прикладі воксельних 3D моделей.
\end{itemize}